\answer{ex:12:1}{$0.60$(상관이 없으면 $0.18$), 극한 $\sqrt{c/(rT)}\approx0.58$: 뉴런 백 개로는 “있다/없다”만 구별한다.}
\answer{ex:12:2}{(a)~스파이크 $10^8$개, $\approx10\,\text{mJ}$: 뇌 전체 $3$~J의 $0.3\,\%$. (b)~$3$~J, $\approx300$배.}
\answer{ex:12:3}{(a)~$\theta_A\in[2,3)$, $\theta_B\in[1,2)$; 다른 도형은 최대 $2$와 $1$을 준다. (b)~$10\cdot96^2\approx9.2\cdot10^4$.}
\answer{ex:12:4}{(a)~$T_d=0.85\,\tau_a=0.34\,\text{s}$. (b)~$T_d\propto\tau_a$이고 $I$와 무관하다: $\tau_a\approx3.5\,\text{s}$(시뮬레이션: $3.1\,\text{s}$).}
\answer{ex:12:5}{$p\approx0.17$. 비율은 천천히 떨어진다: $N=8$, $12$, $24$, $48$, $96$, $192$에서 $0.90$, $0.88$, $0.86$, $0.84$, $0.82$, $0.79$.}
\answer{ex:12:6}{$\tau_{\text{tr}}=20$, $50$, $200\ms$에서 열 번 모두 오류 없음($30\ms$에서는 열 번 중 한 번 실패); $5$--$10\ms$에서는 약 $0.5$, $0.5$--$2\,\text{s}$에서는 품질이 $0.86$까지 떨어진다. 상한: 흔적이 휴지 동안 사라져야 한다, $e^{-0.3\,\text{s}/\tau_{\text{tr}}}\ll1$.}
\answer{ex:12:7}{지연 하나로는 부족하다: 창 $\pm5\ms$로는 $5$에서 $20\ms$까지의 $\Delta x/v$를 덮지 못한다. $5<d_1\le10\ms$와 $15\le d_2\le d_1+10\ms$인 억제 가지 두 개.}
\answer{ex:12:8}{\texttt{two\_stage}: $\Delta=2$, 한 번 뒤집으면 동점, 두 번부터 답이 바뀐다. 1의 계수기: 한 번 뒤집기; $p=0.1$에서 $0.57$.}
